\section{Proof status, archive contents, and reproducibility}
\label{sec:status}

\subsection{The three evidence levels}
The mathematical deliverable is the human argument in this report,
including its numerical bounds and finite-arithmetic interpretation.
The functional-equation foundation is reused as a written proof from
the included upstream archive. The uniform rational-polynomial and
winding arguments supply the finite infrastructure required by the
counting proof; no fixed-degree FTA schema is silently substituted for
the variable-degree construction.

The executable evidence consists of exact rational comparisons for the
worked simple zero and exact integer calculations of the regular-height
bounds. The latter program computes cutoff sizes and verifies their
elementary inequalities; it does not execute the enormous general
height search. A finite run is evidence about that run. The uniform
mathematical proof is the reason the algorithm has the stated behavior
on all inputs.

The Lean deliverable is a concrete interface and a module-by-module
completion plan. No new analytic Lean theorem has been kernel-checked
in this package. The interface has no \code{sorry} proof and adds no
analytic axiom, but no instance of its structure of final obligations
is constructed. A specification without a supplied proof is not a proof. The essential
completion target is an actual term of
\code{Prf PRAAx [] phi}; only then does
\code{pra_iff_native} transfer it to the concise native calculus.

\subsection{Files}
\begin{center}
\begin{tabular}{p{.43\textwidth}p{.49\textwidth}}
\toprule
Path & Purpose\\
\midrule
\code{zeta_zeros_ipra.tex}, \code{sections/}
 & Complete source of this report.\\
\code{zeta_zeros_ipra.pdf}
 & Compiled, visually inspected report.\\
\code{lean/ObjectTargets.lean}
 & Object-language sentence constructors, the verified repository bridge
   as an interface, and explicit final proof obligations.\\
\code{lean/README.md}
 & Proposed build location, dependency versions and completion gates.\\
\code{code/rational_zeta.py}
 & Exact rational zeta enclosures and the simple-zero certificate.\\
\code{code/good_height_bounds.py}
 & Exact integer size calculations for the theoretical regular-height
   construction.\\
\code{verification/}
 & Actual computation outputs and verification notes.\\
\code{sources/}
 & Pinned repository source snapshots and provenance manifest.\\
\path{dependencies/zeta_functional_equation_ipra.zip}
 & Unchanged upstream functional-equation proof, finite contracts and
   explicitly incomplete Lean outline.\\
\code{SHA256SUMS}
 & Hashes for the files in this archive.\\
\bottomrule
\end{tabular}
\end{center}

\subsection{Reproduction}
With Python 3.12 or a compatible later Python, the two scripts require
only the standard library. The main report uses ordinary LaTeX packages.
Run the supplied \code{build.sh} from the archive root to regenerate
the computation outputs and compile the PDF. The exact build used for
this delivery and the observed results are recorded in
\code{verification/STATUS.md}.

The repository snapshot is read-only in this work. Source files copied
from it were checked against their Git blob hashes, and their SHA-256
hashes are recorded. The included functional-equation ZIP has SHA-256
\begin{center}\small\ttfamily
78ddf1de0c4c5fc3e6a553efdeda96e617a646f7a02e62f84481d57c8691e4de.
\end{center}
No claim about newer repository revisions is needed for the argument.
